% New exposition and proofs. No source-paper text is reproduced.
\section{Dirichlet constituents and zeros of umbral shadows}
\label{sec:shadow-dirichlet}
This section identifies exactly which scalar functions obtained from an umbral
shadow inherit the Dirichlet Riemann hypothesis.  It also determines the zeros
of every component involving two distinct primitive Dirichlet characters by
applying the theorem of Saias and Weingartner \cite[pp.~335--336]{SW2009}.
In particular, the first component for $A_5^4D_4$ has infinitely many zeros to
the right of its expected critical line, and even in its half-plane of absolute
convergence.  This last conclusion is unconditional.

Put $q=e^{2\pi i\tau}$ and, for $m\geq2$ and $r\in\mathbb Z/2m\mathbb Z$, put
\begin{equation}\label{eq:unary-definition}
 S_{m,r}(\tau)=\sum_{\substack{n\in\mathbb Z\\n\equiv r\pmod{2m}}}
 n q^{n^2/(4m)}.
\end{equation}
Thus $S_{m,-r}=-S_{m,r}$ and $S_{m,0}=S_{m,m}=0$.
This is the normalization of Cheng, Duncan and Harvey
\cite[Sect.~3.4, equation (3.39), p.~37]{CDH2014}; in particular, its exponent denominator
is $4m$.  For a series
\[
 H(\tau)=\sum_{n\geq1}c(n)nq^{n^2/(4m)}
\]
with periodic coefficients we use the Dirichlet transform
\begin{equation}\label{eq:shadow-L-normalization}
 L(H,s)=\sum_{n\geq1}c(n)n\left(\frac{n^2}{4m}\right)^{-s}
       =(4m)^sD_c(2s-1),\qquad
 D_c(z)=\sum_{n\geq1}\frac{c(n)}{n^z}.
\end{equation}
The defining sums converge absolutely for $\Re s>1$ and $\Re z>1$,
respectively.  They thereafter mean their meromorphic continuations.

\begin{theorem}[Character constituents at every lambency]\label{thm:all-shadow-characters}
Let $N=2m$, let $d\mid N$, put $Q=N/d$, and let $\chi$ be an odd Dirichlet
character modulo $Q$.  Define
\[
 G_{d,\chi}(\tau)=\sum_{r\bmod Q}\chi(r)S_{m,dr}(\tau).
\]
Then
\begin{equation}\label{eq:all-shadow-character-L}
 L(G_{d,\chi},s)=2d\left(\frac{4m}{d^2}\right)^sL(2s-1,\chi).
\end{equation}
If $\chi^*$ is the primitive character inducing $\chi$, of conductor $f$, then
\begin{equation}\label{eq:imprimitive-shadow}
 L(2s-1,\chi)=L(2s-1,\chi^*)
 \prod_{\substack{p\mid Q\\p\nmid f}}
       \left(1-\chi^*(p)p^{1-2s}\right).
\end{equation}
Consequently the Dirichlet Riemann hypothesis implies that every zero of
this nonzero constituent in $1/2<\Re s<1$ satisfies $\Re s=3/4$.
The omitted Euler factors have their zeros on $\Re s=1/2$.
Every component of every fixed umbral shadow or twined shadow is a finite
linear combination of the $G_{d,\chi}$.
\end{theorem}
\begin{proof}
In the sum defining $G_{d,\chi}$, set $n=dk$.  Summing over residue classes
first gives
\[d\sum_{k\in\mathbb Z}\chi(k)kq^{d^2k^2/(4m)}.\]
Since $\chi(-1)=-1$, the positive and negative terms agree.  Mellin expansion
in the absolutely convergent half-plane proves
\eqref{eq:all-shadow-character-L}.  The usual Euler product proves
\eqref{eq:imprimitive-shadow} there and hence everywhere by continuation.
For a prime in that product $|\chi^*(p)|=1$, so a zero of its factor has
$p^{1-2\Re s}=1$.  The claimed zero lines follow.

For the spanning assertion, write a component as
$H=\sum_{r\bmod N}h_rS_{m,r}$.  Its coefficient in
\eqref{eq:shadow-L-normalization} is the odd periodic function
$c(n)=h_n-h_{-n}$.  On the stratum $\gcd(n,N)=d$, the function $k\mapsto c(dk)$
is an odd function on $(\mathbb Z/Q\mathbb Z)^\times$.  Orthogonality of
characters gives
\begin{equation}\label{eq:character-stratum-coefficients}
 c(dk)=\sum_{\chi\bmod Q}C_{d,\chi}\chi(k),\qquad
 C_{d,\chi}=\frac1{\varphi(Q)}
    \sum_{u\in(\mathbb Z/Q\mathbb Z)^\times}c(du)\overline{\chi(u)}.
\end{equation}
Changing $u$ to $-u$ shows that $C_{d,\chi}=0$ for even $\chi$.
Thus $H=\frac12\sum_{d,\chi}C_{d,\chi}G_{d,\chi}$.
The umbral and twined shadows are such linear combinations by their defining
matrix formula $S_g^X=\Omega_g^XS_m$
\cite[Sect.~4.1 and Sect.~5.1]{CDH2014}.  No restriction on the Coxeter number or on
$g$ was used.
\end{proof}

The distinction between a constituent and a sum of constituents is essential.
In particular, a primitive character may occur on several gcd strata, and
adding those strata produces a Dirichlet polynomial which need not be a
product of omitted Euler factors.

\begin{theorem}[Classification of zeros of a scalar shadow]\label{thm:shadow-zero-classification}
Let $H$ be a nonzero fixed linear combination of the $S_{m,r}$, and write its
coefficient function $c$ as above.  There is a unique finite expansion
\begin{equation}\label{eq:primitive-shadow-expansion}
 D_c(z)=\sum_{\psi\ {
m primitive}}P_\psi(z)L(z,\psi)
\end{equation}
in which each $P_\psi$ is a Dirichlet polynomial.  All characters occurring
are odd.

If at least two $P_\psi$ are nonzero, there is an $\epsilon>0$ such that,
for every $3/4<\alpha<\beta\leq1+\epsilon$, the number of distinct zeros of
$L(H,s)$ in
\[
 \alpha<\Re s<\beta,\qquad |\Im s|\leq T
\]
is bounded above and below by positive constant multiples of $T$, for all
sufficiently large $T$.  In particular there are infinitely many zeros in
$\Re s>1$.  These conclusions need no Riemann hypothesis.

If exactly one $P_\psi$ is nonzero, the zeros are those of
$P_\psi(2s-1)$ and $L(2s-1,\psi)$, with multiplicity, apart from no further
factors.  The Dirichlet Riemann hypothesis controls the second factor only.
\end{theorem}
\begin{proof}
Split the sum for $D_c$ by $\gcd(n,N)=d$ and use
\eqref{eq:character-stratum-coefficients}.  This yields
\[
 D_c(z)=\sum_{d\mid N}\sum_{\chi\bmod N/d}
 C_{d,\chi}d^{-z}L(z,\chi).
\]
Replacing each imprimitive $L$-function by its primitive one times its finite
Euler product gives \eqref{eq:primitive-shadow-expansion}.
Uniqueness is precisely the Primitive Dirichlet Character Basis theorem of
Saias and Weingartner \cite[p.~335]{SW2009}.  Thus having two nonzero
primitive terms is equivalent to $D_c$ not being a Dirichlet polynomial times
one Dirichlet $L$-function.  Their main theorem
\cite[pp.~335--336]{SW2009} supplies $\eta>0$ and upper and lower linear
zero-counting bounds in every strip
$1/2<\sigma_1<\sigma_2\leq1+\eta$ for $D_c$.
Set $\sigma_1=2\alpha-1$, $\sigma_2=2\beta-1$ and
$\epsilon=\eta/2$.  The affine change $z=2s-1$ doubles the height and
preserves multiplicity; the factor $(4m)^s$ has no zero.  This proves the
assertion.  In the one-character case $\psi$ is odd and hence nonprincipal,
so its $L$-function is entire.  The stated product therefore has no pole
which could cancel a zero.
\end{proof}

This theorem answers the question for every lambency at once whenever a
component mixes different primitive conductors: it has zeros off
$\Re s=3/4$ and infinitely many zeros in $\Re s>1$.  The hypothesis refers
to the final primitive expansion, so terms which cancel must first be
combined.  Merely observing nonmultiplicative coefficients is insufficient,
since a single Dirichlet $L$-function multiplied by a Dirichlet polynomial
can also have nonmultiplicative coefficients.

\subsection{The full lambency-six calculation}
For $m=6$ put
\[
 A=S_{6,1}+S_{6,5},\qquad B=S_{6,1}-S_{6,5},\qquad z=2s-1.
\]
Taking positive and negative integers in their residue classes gives
\begin{align}
 A(\tau)&=\sum_{\substack{n\geq1\\(n,6)=1}}\chi_{-4}(n)nq^{n^2/24},&
 L(A,s)&=24^s(1+3^{-z})L(z,\chi_{-4}),\label{eq:l6-A}\\
 B(\tau)&=\sum_{\substack{n\geq1\\(n,6)=1}}\chi_{-3}(n)nq^{n^2/24},&
 L(B,s)&=24^s(1+2^{-z})L(z,\chi_{-3}).\label{eq:l6-B}
\end{align}
The remaining components are
\begin{align}
 S_{6,2}(\tau)&=2\sum_{\substack{k\geq1\\(k,6)=1}}
 \chi_{-3}(k)kq^{k^2/6},&
 L(S_{6,2},s)&=2\,6^s(1+2^{-z})L(z,\chi_{-3}),\label{eq:l6-2}\\
 S_{6,3}(\tau)&=3\sum_{k\geq1}\chi_{-4}(k)kq^{3k^2/8},&
 L(S_{6,3},s)&=3(8/3)^sL(z,\chi_{-4}),\label{eq:l6-3}\\
 S_{6,4}(\tau)&=4\sum_{k\geq1}\chi_{-3}(k)kq^{2k^2/3},&
 L(S_{6,4},s)&=4(3/2)^sL(z,\chi_{-3}).\label{eq:l6-4}
\end{align}
For example, in $S_{6,2}$ write $n=2k$ and distinguish $k\equiv1,5\pmod6$;
the second residue class comes from a negative $n$.  The other identities
follow in exactly the same way, by the displayed congruence classes.

\begin{theorem}[The first mixed umbral component]\label{thm:l6-mixed-zeros}
For $X=A_5^4D_4$, the first untwined shadow component has transform
\begin{equation}\label{eq:l6-mixed-L}
 L(S^X_1,s)=24^s\left\{
 3(1+3^{1-2s})L(2s-1,\chi_{-4})+
 2(1+2^{1-2s})L(2s-1,\chi_{-3})\right\}.
\end{equation}
It has infinitely many zeros off $\Re s=3/4$, including infinitely many
zeros in $\Re s>1$, with the stripwise linear bounds of
Theorem~\ref{thm:shadow-zero-classification}.
The same conclusions hold for $S^X_5$.
\end{theorem}
\begin{proof}
For this root system, the defining ADE matrices give
$S^X_1=5S_{6,1}+S_{6,5}=3A+2B$ and
$S^X_5=S_{6,1}+5S_{6,5}=3A-2B$.
Substitution of \eqref{eq:l6-A}--\eqref{eq:l6-B} proves the transforms.
Both primitive terms are nonzero, and their conductors are $4$ and $3$,
so Theorem~\ref{thm:shadow-zero-classification} applies.
Alternatively, for the first component its periodic coefficients on
$1,5,7,11\pmod{12}$ are $5,1,-1,-5$.
If $D_c=P L(\cdot,\chi)$ with $P=\sum_{d\leq M}b_dd^{-z}$, then coefficient
uniqueness gives $b_1=5$ and $c(p)=5\chi(p)$ at every prime $p>M$ not dividing
the conductor of $\chi$.  There are arbitrarily large primes
$p\equiv5\pmod{12}$, and there $|c(p)|=1$, a contradiction.
\end{proof}

\subsection{Every lambency-six twining and its symmetric powers}
Let $G=\mathrm{GL}_2(\mathbb F_3)$.  Write $a(g)$ for the number of fixed
lines in $\mathbb F_3^2$, and $b(g)$ for the number of those lines on which
$g$ acts by $+1$ minus the number on which it acts by $-1$.  On the three
partitions of these four lines into two unordered pairs, put
$c(g)=\#\mathrm{Fix}(g)-1$, and let $\delta(g)$ be the sign of this
permutation.  Equivalently $\delta(g)=\det_{\mathbb F_3}(g)\in\{1,-1\}$.
The shadow definition of Cheng, Duncan and Harvey
\cite[Sect.~5.1, equations (5.1)--(5.6), and Table~28, p.~95]{CDH2014} gives
\begin{equation}\label{eq:l6-full-shadow}
 S^X_g=
 \left( (a+1)S_{6,1}+\delta S_{6,5},\;
 bS_{6,2},\;(a+c)S_{6,3},\;bS_{6,4},\;
 \delta S_{6,1}+(a+1)S_{6,5}\right).
\end{equation}
One may check this directly: the four $A_5$ diagrams give $a$ on odd
indices and $b$ on even indices; the $D_4$ triality diagram gives $1$ on
indices $1,5$, $c$ on index $3$, and the two off-diagonal entries
$\delta$.  Its two entries on the even indices cancel because
$S_{6,-r}=-S_{6,r}$.

\begin{theorem}[Symmetric powers in the twined shadow]\label{thm:l6-sympowers}
Choose the faithful two-dimensional complex representation $V$ of $G$
whose trace on one of the order-eight classes is $\sqrt{-2}$.  Then
\[
 \det V=\delta,\qquad
 a-1=\operatorname{tr}\operatorname{Sym}^2V,\qquad
 b=\operatorname{tr}\operatorname{Sym}^3V,\qquad
 a+c=1+\operatorname{tr}\operatorname{Sym}^4V.
\]
Moreover
\[
 V\otimes\operatorname{Sym}^2V
   \simeq\operatorname{Sym}^3V\oplus(V\otimes\det V),\qquad
 V\otimes\det V\simeq\overline V.
\]
The centre $-I$ acts on $\operatorname{Sym}^kV$ by $(-1)^k$.
The nonzero isotypic combinations of every fixed twining are constant
multiples of the five Dirichlet constituents in
\eqref{eq:l6-A}--\eqref{eq:l6-4}.
\end{theorem}
\begin{proof}
Here is the complete class calculation; the two order-eight classes are
combined only in the columns which have equal values.  The faithful
character is the two-dimensional character in the character table of
$\mathrm{GL}_2(3)$ \cite[Table~15, p.~91]{CDH2014}.
\begin{center}
\begin{tabular}{c|rrrrrrr}
class&$1A$&$2A$&$2B$&$4A$&$3A$&$6A$&$8A,8B$\\\hline
size&1&1&12&6&8&8&6,6\\
$\operatorname{tr}V$&2&$-2$&0&0&$-1$&1&$\pm\sqrt{-2}$\\
$\delta$&1&1&$-1$&1&1&1&$-1$\\
$a$&4&4&2&0&1&1&0\\
$b$&4&$-4$&0&0&1&$-1$&0\\
$c$&2&2&0&2&$-1$&$-1$&0
\end{tabular}
\end{center}
The last three rows can also be obtained without a character table by acting
on the four lines and their three pair partitions.  For trace $t$ and
determinant $\delta$ the symmetric-power traces are
\[
 1,\quad t,\quad t^2-\delta,\quad t^3-2\delta t,
 \quad t^4-3\delta t^2+\delta^2.
\]
Substitution in the displayed table proves every character equality.
For completeness, the tensor identity has a canonical proof independent of
this table.  Symmetrization maps $V\otimes\operatorname{Sym}^2V$ onto
$\operatorname{Sym}^3V$; its kernel is the copy of
$V\otimes\bigwedge^2V$ generated by
$v\otimes(uw)-u\otimes(vw)$.  Complete reducibility gives the direct sum.
Multiplication of the trace row by $\delta$ gives its complex conjugate,
proving the second isomorphism.  Faithfulness implies that the nontrivial
central involution acts by $-1$ on $V$, and hence by $(-1)^k$ on its
$k$th symmetric power.
Finally \eqref{eq:l6-full-shadow} gives
\begin{align*}
 S^X_{g,1}+S^X_{g,5}&=(a+1+\delta)A,\\
 S^X_{g,1}-S^X_{g,5}&=(a+1-\delta)B,
\end{align*}
while the other three entries are $bS_{6,2}$, $(a+c)S_{6,3}$, and
$bS_{6,4}$.  These are exactly the asserted constituents.
\end{proof}

The zero cases must be excluded when making any Riemann-hypothesis
statement.  Explicitly the respective multipliers of
$(A,B,S_{6,2},S_{6,3},S_{6,4})$, where the first two mean the sum and
difference of the twined components, are
\begin{center}
\begin{tabular}{c|rrrrr}
class&$A$&$B$&$S_{6,2}$&$S_{6,3}$&$S_{6,4}$\\\hline
$1A$&6&4&4&6&4\\
$2A$&6&4&$-4$&6&$-4$\\
$2B$&2&4&0&2&0\\
$4A$&2&0&0&2&0\\
$3A$&3&1&1&0&1\\
$6A$&3&1&$-1$&0&$-1$\\
$8A,8B$&0&2&0&0&0
\end{tabular}
\end{center}
Every nonzero entry inherits the zero line of its constituent.  A zero entry
is the identically zero function and has no discrete zero-counting problem.
For a single first component, as opposed to its sum and difference with the
fifth, the coefficients of $A$ and $B$ are
$(a+1+\delta)/2$ and $(a+1-\delta)/2$.
Thus the first and fifth components have the unconditional mixed-character
zero phenomenon on $1A,2A,2B,3A,6A$; they are single-character constituents
on $4A,8A,8B$.

If an elliptic curve $E/\mathbb Q$ has full mod-$3$ image, the Weil pairing
identifies $\det\overline\rho_{E,3}$ with the mod-$3$ cyclotomic character.
Consequently, at an unramified Frobenius,
$\delta(\overline\rho_{E,3}(\mathrm{Frob}_p))=\chi_{-3}(p)$.
The remaining character rows are traces of the finite-image Artin
representations $\operatorname{Sym}^j(V\circ\overline\rho_{E,3})$.
Their Artin $L$-functions are different objects from the unary transforms
of a fixed twining; the Dirichlet theorem directly controls the determinant
character but supplies no Riemann hypothesis for these higher-dimensional
Artin representations.

\subsection{The elementary lattice factorizations}
The following formulas use $(v,v)/2$ as the quadratic form for $A_2,D_4,E_8$,
and the ordinary Gaussian norm for $\mathbb Z[i]$:
\begin{align}
 Z_{A_2}(s)&=6\zeta(s)L(s,\chi_{-3}),&
 Z_{\mathbb Z[i]}(s)&=4\zeta(s)L(s,\chi_{-4}),\label{eq:CM-Epstein}\\
 Z_{D_4}(s)&=24(1-2^{1-s})\zeta(s)\zeta(s-1),&
 Z_{E_8}(s)&=240\zeta(s)\zeta(s-3).\label{eq:DE-Epstein}
\end{align}
For the first pair, $\mathbb Z[\omega]$ and $\mathbb Z[i]$ are principal
ideal domains with respectively six and four units.  Grouping nonzero
elements by their principal ideals gives the respective Dedekind zeta
functions, which factor as displayed.  For $D_4$, use
$D_4=\{x\in\mathbb Z^4:\sum x_i\equiv0\pmod2\}$.
The condition $\sum x_i^2=2n$ already imposes this parity condition, so
Jacobi's four-square formula gives
\[
 r_{D_4}(n)=r_4(2n)=8\sum_{d\mid2n,\,4\nmid d}d
       =24\sum_{d\mid n,\,d\ {
m odd}}d.
\]
Taking its Dirichlet series proves the third identity.  The theta series of
$E_8$ is the unique normalized weight-four modular form on
$\mathrm{SL}_2(\mathbb Z)$, namely $E_4=1+240\sum_{n\geq1}\sigma_3(n)q^n$;
this proves the fourth.

Thus the Dirichlet Riemann hypothesis places the nontrivial Dirichlet-factor
zeros on $\Re s=1/2$ for the first pair, on $\Re s=1/2,3/2$ for $D_4$,
and on $\Re s=1/2,7/2$ for $E_8$.  The factor $1-2^{1-s}$ contributes
zeros on $\Re s=1$, except that its zero at $s=1$ cancels the pole of
$\zeta(s)$ and is not a zero of $Z_{D_4}$.
Similarly, the signed transform
\[
 (4^s-2)\,6\zeta(s)L(s,\chi_{-3})
\]
has all its zeros in $0<\Re s<1$ on $\Re s=1/2$, since the additional
factor vanishes precisely at
$s=\tfrac12+\pi i k/\log2$, $k\in\mathbb Z$.
At lambency two, Jacobi's identity
$\eta(\tau)^3=\sum_{n\geq1}\chi_{-4}(n)nq^{n^2/8}$ gives
$L(24\eta^3,s)=24\,8^sL(2s-1,\chi_{-4})$.
